\documentclass[12pt,a4paper]{article}
\usepackage[utf8]{inputenc}
\usepackage[T1]{fontenc}
\usepackage[polish]{babel}
\usepackage{polski}
\usepackage{geometry}
\usepackage{booktabs}
\usepackage{tabularx}
\usepackage{parskip}
\usepackage{titlesec}

\geometry{margin=2.5cm}

\title{\textbf{Systemy Wbudowane Lista 3}}

\begin{document}

\maketitle

\section{Zadanie 1 -- Diagram komponentów}

\textbf{Jednostka Główna} to węzeł nadrzędny realizujący logikę kolejkowania zadań i wyznaczania trasy kabiny. Pamięć konfiguracyjna i zegar RTC podłączone są przez I2C, moduł komunikacji zewnętrznej przez UART.

\textbf{Kabina} zawiera sterownik Slave obsługujący wszystkie podzespoły. Przyciski i czujniki binarne podłączone są przez GPIO, wyświetlacz i czytnik kart przez SPI, czujnik wagi przez ADC, głośnik przez I2S, falownik silnika i siłownik drzwi przez PWM.

\textbf{Węzeł Piętrowy} obsługuje przyciski i wyświetlacz na każdym piętrze. Komunikuje się z Jednostką Główną przez magistralę CAN, podobnie jak Kabina. CAN został wybrany ze względu na odporność na zakłócenia elektromagnetyczne panujące w szybie windy.

\textbf{Systemy Zewnętrzne} -- PPOŻ i panel serwisowy -- komunikują się przez moduł zewnętrzny oddzielony od magistrali CAN, co zapewnia niezależność systemu bezpieczeństwa.

\section{Zadanie 2 -- Dane w systemie}

\subsection{Pozycja kabiny}

\begin{tabularx}{\textwidth}{lX}
\toprule
\textbf{Atrybut} & \textbf{Wartość} \\
\midrule
Wielkość fizyczna & Położenie liniowe \\
Jednostka & Milimetry [mm] \\
Zakres & 0 -- maksymalna wysokość szybu \\
Rozdzielczość przechowywana & 1 mm \\
Rozdzielczość wyświetlana & Numer piętra (liczba całkowita) \\
Format wymiany & \texttt{uint32} -- surowa wartość z enkodera \\
Źródło & Czujnik pozycji \\
Odbiorcy & Master, wyświetlacz kabiny, wyświetlacz piętrowy \\
\bottomrule
\end{tabularx}

\subsection{Waga kabiny}

\begin{tabularx}{\textwidth}{lX}
\toprule
\textbf{Atrybut} & \textbf{Wartość} \\
\midrule
Wielkość fizyczna & Masa \\
Jednostka & Kilogramy [kg] \\
Zakres & 0 -- 150\% udźwigu nominalnego \\
Rozdzielczość przechowywana & 0,1 kg \\
Rozdzielczość wyświetlana & Niewyświetlana użytkownikowi \\
Format wymiany & \texttt{float32} -- odczyt ADC przeliczony na kg \\
Źródło & Czujnik wagi \\
Odbiorcy & Sterownik kabiny, Master \\
\bottomrule
\end{tabularx}

\subsection{Żądania jazdy}

\begin{tabularx}{\textwidth}{lX}
\toprule
\textbf{Atrybut} & \textbf{Wartość} \\
\midrule
Wielkość fizyczna & Brak -- dane logiczne \\
Jednostka & Brak \\
Zakres & Piętro: 0 -- N, kierunek: \{góra, dół\} \\
Rozdzielczość przechowywana & 1 piętro \\
Rozdzielczość wyświetlana & Podświetlenie przycisku wybranego piętra \\
Format wymiany & \texttt{\{floor: uint8, direction: enum, source: enum\}} \\
Źródło & Panel przycisków kabiny, przyciski piętrowe \\
Odbiorcy & Master -- kolejka zadań \\
\bottomrule
\end{tabularx}

\subsection{Stan drzwi}

\begin{tabularx}{\textwidth}{lX}
\toprule
\textbf{Atrybut} & \textbf{Wartość} \\
\midrule
Wielkość fizyczna & Brak -- dane logiczne \\
Jednostka & Brak \\
Zakres & \{otwarte, zamknięte, w ruchu, zablokowane\} \\
Rozdzielczość przechowywana & Enum 2 bity \\
Rozdzielczość wyświetlana & Niewyświetlany bezpośrednio \\
Format wymiany & \texttt{uint8} -- enum stanu \\
Źródło & Fotokomórka drzwi, siłownik drzwi \\
Odbiorcy & Sterownik kabiny, Master \\
\bottomrule
\end{tabularx}

\subsection{Prędkość kabiny}

\begin{tabularx}{\textwidth}{lX}
\toprule
\textbf{Atrybut} & \textbf{Wartość} \\
\midrule
Wielkość fizyczna & Prędkość liniowa \\
Jednostka & Milimetry na sekundę [mm/s] \\
Zakres & 0 -- 2500 mm/s \\
Rozdzielczość przechowywana & 1 mm/s \\
Rozdzielczość wyświetlana & Niewyświetlana użytkownikowi \\
Format wymiany & \texttt{uint16} \\
Źródło & Wyliczana z czujnika pozycji \\
Odbiorcy & Sterownik kabiny, falownik silnika \\
\bottomrule
\end{tabularx}

\subsection{Tryb pracy systemu}

\begin{tabularx}{\textwidth}{lX}
\toprule
\textbf{Atrybut} & \textbf{Wartość} \\
\midrule
Wielkość fizyczna & Brak -- dane logiczne \\
Jednostka & Brak \\
Zakres & \{normalny, serwisowy, awaryjny PPOŻ, przeciążenie\} \\
Rozdzielczość przechowywana & Enum 3 bity \\
Rozdzielczość wyświetlana & Komunikat na wyświetlaczu i głośniku \\
Format wymiany & \texttt{uint8} -- enum trybu \\
Źródło & Master, System PPOŻ, Panel serwisowy \\
Odbiorcy & Wszystkie komponenty systemu \\
\bottomrule
\end{tabularx}

\subsection{Parametry konfiguracyjne}

\begin{tabularx}{\textwidth}{lX}
\toprule
\textbf{Atrybut} & \textbf{Wartość} \\
\midrule
Wielkość fizyczna & Brak -- nastawy serwisowe \\
Jednostka & Zależna od parametru \\
Zakres & Czas otwarcia drzwi: 1--10 s, prędkość nominalna: 100--2500 mm/s, piętro bazowe: 0--N \\
Rozdzielczość przechowywana & Zależna od parametru \\
Rozdzielczość wyświetlana & Wartość liczbowa na panelu serwisowym \\
Format wymiany & \texttt{\{param\_id: uint8, value: uint32\}} \\
Źródło & Panel serwisowy \\
Odbiorcy & Pamięć EEPROM, Master \\
\bottomrule
\end{tabularx}

\subsection{Uprawnienia dostępu}

\begin{tabularx}{\textwidth}{lX}
\toprule
\textbf{Atrybut} & \textbf{Wartość} \\
\midrule
Wielkość fizyczna & Brak -- dane logiczne \\
Jednostka & Brak \\
Zakres & ID karty: 0 -- $2^{32}$, poziom dostępu: \{pasażer, serwis, służby\} \\
Rozdzielczość przechowywana & 32-bitowe ID karty \\
Rozdzielczość wyświetlana & Niewyświetlane \\
Format wymiany & \texttt{\{card\_id: uint32, access\_level: uint8\}} \\
Źródło & Czytnik kart dostępu \\
Odbiorcy & Master, pamięć konfiguracji \\
\bottomrule
\end{tabularx}

\subsection{Relacje między danymi}

\textbf{Żądania jazdy} trafiają do Mastera, który porównuje je z \textbf{pozycją kabiny} i generuje polecenia ruchu.

\textbf{Waga kabiny} nadpisuje żądania jazdy -- przekroczenie progu wymusza tryb przeciążenia i blokuje ruch.

\textbf{Stan drzwi} blokuje zmianę \textbf{prędkości kabiny} -- jazda możliwa tylko przy drzwiach zamkniętych i zaryglowanych.

\textbf{Tryb pracy} jest nadrzędny dla całego systemu -- tryb awaryjny PPOŻ anuluje wszystkie żądania i wymusza zjazd na piętro ewakuacyjne.

\textbf{Uprawnienia dostępu} filtrują żądania jazdy na piętra z ograniczonym dostępem.

\textbf{Parametry konfiguracyjne} z EEPROM definiują wartości graniczne dla prędkości, czasu otwarcia drzwi i piętra bazowego.

\end{document}